% 第五章：Linux 内核实现

\section{Linux 内核实现} \label{sec:kernel}

\subsection{数据结构}

内核侧的实体关系是"任务—作业—DAG—池—表"五层：

\begin{itemize}
  \item \texttt{struct task\_struct} 内嵌 \texttt{struct sched\_yat\_casched\_entity}，保存 \texttt{job} 指针、\texttt{run\_list}、\texttt{last\_cpu}、\texttt{vruntime}。
  \item \texttt{struct yat\_job} 是调度单元：\texttt{job\_id}、\texttt{job\_priority}、\texttt{wcet}、\texttt{state}（WAITING/READY/RUNNING/COMPLETED）、\texttt{pending\_predecessors}、\texttt{assigned\_cpu}、\texttt{total\_exec}，以及 \texttt{pred/succ} 两条链表。
  \item \texttt{struct yat\_dag\_task} 是一次 DAG 实例：\texttt{dag\_id}、\texttt{total\_jobs}、\texttt{completed\_jobs}、\texttt{jobs} 链表、自旋锁。
  \item \texttt{struct yat\_ready\_job\_pool} 是就绪池：按 \texttt{job\_priority} 插入的双向链表 + \texttt{pool->lock}。
  \item \texttt{struct yat\_cpu\_history\_table} 每 CPU 一条按时间递增的历史链表（上限 512 条）；\texttt{struct yat\_accel\_table} 是按 \texttt{(pid,cpu)} 哈希的加速表，只作为可观测性缓存，不参与控制流。
\end{itemize}

\subsection{作业生命周期}

早期实现每次唤醒都新建一个 \texttt{yat\_job}，导致"每唤醒泄漏一个作业、\texttt{dag->total\_jobs} 无界增长、历史条目与新 \texttt{job\_id} 对不上"。现在作业与任务同寿：首次入队时创建，之后只更新状态，销毁点唯一——\texttt{task\_dead}。这条路径是 \texttt{finish\_task\_switch} 里调的，任务已完成最后一次切换，\texttt{task\_struct} 仍有效，因此可以安全地取指针并清空：

\begin{verbatim}
static void task_dead_yat_casched(struct task_struct *p)
{
    struct yat_job *job = READ_ONCE(p->yat_casched.job);
    unsigned long flags;
    if (!job) return;
    WRITE_ONCE(p->yat_casched.job, NULL);
    local_irq_save(flags);
    yat_destroy_job(job);
    local_irq_restore(flags);
}
\end{verbatim}

\texttt{yat\_destroy\_job()} 的顺序是先摘 \texttt{pool->lock} 下的链表节点，再摘 \texttt{dag->lock} 下的 \texttt{job\_list}，最后在锁外 \texttt{kfree}——这与文件头登记的锁序一致，也不需要额外加锁。

\subsection{唤醒与选核}

\texttt{select\_task\_rq\_yat\_casched()} 是热路径，分三段：

\begin{enumerate}
  \item \textbf{窗口复用}。读 \texttt{job->assigned\_cpu}，若目标核仍在线、仍允许、且在 \texttt{YAT\_DECISION\_MIN\_INTERVAL\_NS} 内，分两种情况直接返回：目标核空闲（\texttt{select\_sticky\_idle}）、或间隔未到（\texttt{select\_sticky\_intv}）。这两条计数在实测中分别是 195 与 0，说明"窗口内复用"确实拦下了约四成选核。
  \item \textbf{亲和性过滤}。候选核必须满足 \texttt{cpu\_active(cpu)} 且 \texttt{cpumask\_test\_cpu(cpu, p->cpus\_ptr)}。这一条是本轮修的：原实现把 \texttt{is\_cpu\_allowed()} 的兜底当成"不用校验"，结果 \texttt{best\_core} 可能落在任务不被允许的核上，进而污染历史表与自助计数。
  \item \textbf{全量选择}。\texttt{yat\_find\_best\_core\_allowed()} 遍历在线核，算 \texttt{recency}、CRP 比率、benefit，取最大；并列时用 \texttt{Impact} 破平。找不到核时返回 \texttt{task\_cpu}，由通用层兜底。
\end{enumerate}

\texttt{set\_cpus\_allowed\_yat\_casched()} 直接调 \texttt{set\_cpus\_allowed\_common()}：它会同时更新 \texttt{cpus\_mask}/\texttt{nr\_cpus\_allowed} 并处理 migrate-disable，少了这一步，\texttt{p->cpus\_ptr} 与 \texttt{cpus\_mask} 会不一致。

\subsection{一次被实证的锁反序及其修复}

镜像 \#18 的串口日志里出现过一条 \texttt{lockdep} 报告，关键行摘录如下：

\begin{verbatim}
WARNING: possible irq lock inversion dependency detected
 swapper/1/0 just changed the state of lock:
 ffff888001090028 (&pool->lock#3){+...}-{2:2}, at: yat_destroy_job+0x36/0x3e0
 but this lock was taken by another, HARDIRQ-safe lock in the past:
 (&rq->__lock)
 Possible interrupt unsafe locking scenario:
       CPU0                    CPU1
       lock(&pool->lock#3);
                               local_irq_disable();
                               lock(&rq->__lock);
                               lock(&pool->lock#3);
       <Interrupt>
       lock(&rq->__lock);
\end{verbatim}

成因清楚：\texttt{pool->lock} 在 balance 路径里总是"持 \texttt{rq\_lock}、关中断"获取，因此是 HARDIRQ-safe 锁；而 \texttt{task\_dead} 在中断开启的上下文里获取它，恰好构成"任务持 \texttt{pool->lock} 时被中断打断、中断里取 \texttt{rq\_lock}"的环。第一次修复落在 debugfs 的 \texttt{show()} 上（改成 \texttt{spin\_lock\_irqsave}），只堵住了一条路径；真正的根因是 \texttt{task\_dead}。

修复方式是把整段销毁放进关中断区间（见上节代码）：临界区内只有链表操作与 \texttt{kfree}，不会睡眠，时延是常数级。修复后重建镜像并复跑两组对照，串口日志中的同类告警计数为 0（第\ref{sec:test}节）。

\subsection{可观测性}

调度器自带 \texttt{debugfs} 接口与 per-CPU 计数器，二者分工明确：计数器是"便宜、可累加、可跨轮次对比"的量，tracepoint 是"带参数、可离线分析"的量。计数器共 31 项（\texttt{debugfs:/sys/kernel/debug/yat\_casched/stats}），关键的有：

\begin{itemize}
  \item \texttt{select\_calls / select\_decided / select\_sticky\_idle / select\_nocore}：选核回调是否被调用、做出了决策、还是走了复用/失败分支；
  \item \texttt{enq\_calls / enq\_cpus\_allowed\_sum}：后者除以前者即"平均可运行核数"，是判断"回调没被调用是因为亲和性"的唯一依据；
  \item \texttt{jobs\_placed / jobs\_nocore}：就绪池是否真的投出了作业；
  \item \texttt{place\_hit / place\_miss / place\_pending}：把"决策的核"与"实际运行的核"在 pick 时刻比对，得到接管率。
\end{itemize}

接管率的定义要写死：$\text{hit}/(\text{hit}+\text{miss}+\text{pending})$。\texttt{placements} 文件是实时快照，任务退出后即空，因此它不能作为留存证据；能留存的是三个累计值。

这 31 个计数器不是仪表盘，而是一组\textbf{可证伪的断言}：本文每一条"系统在工作"的结论，都指定了它会看哪个读数、以及断言为假时读数会怎么变（\autoref{tab:assert}）。这样做的直接好处是排查不靠猜——第\ref{sec:test}节的"空转"问题就是从 \texttt{balance\_runs} 与 \texttt{ttwu\_ok} 的比值读出来的。
ef{sec:test}节的"空转"问题就是从 	exttt{balance\_runs} 与 	exttt{ttwu\_ok} 的比值读出来的。

\begin{table}[H]
\centering
\caption{三条关键断言及其可证伪读数}
\label{tab:assert}
\small
\begin{tabular}{@{}>{\raggedright\arraybackslash}p{3.4cm}>{\raggedright\arraybackslash}p{4.8cm}>{\raggedright\arraybackslash}p{5.4cm}@{}}
\toprule
断言 & 支持它的读数 & 若断言为假会看到什么 \\
\midrule
回调确实被调用 & \texttt{select\_calls}=467 / \texttt{ttwu\_ok}=512 & \texttt{select\_calls} 接近 0 而 \texttt{enq\_calls} 正常 \\
决策确实落地 & 接管率 $906/924=98.1\%$，\texttt{place\_miss}=0，\texttt{pick\_null\_desync}=0 & \texttt{place\_hit} 下降、\texttt{pending} 上升、\texttt{desync} 非 0 \\
优化确实生效 & \texttt{balance\_runs} $757\to9$，\texttt{balance\_noop}=621 & \texttt{balance\_runs} 不降，或 \texttt{balance\_noop}=0 \\
\bottomrule
\end{tabular}
\end{table}

\subsection{亲和性语义与测试}

\texttt{--affinity} 三档必须分清，否则会得出完全错误的结论：

\begin{itemize}
  \item \texttt{shared}：全部任务 pin 到同一核，\texttt{nr\_cpus\_allowed=1}，\texttt{core.c} 直接返回该核，\textbf{回调一次都不会被调用}。这类数据只能谈时间片，不能谈 AJLR。
  \item \texttt{none}：不设亲和性，可跨核；配合 \texttt{--sleep-us} 制造 \texttt{WF\_TTWU}，才是 AJLR 的真实场景。
  \item \texttt{spread}：第 $i$ 个任务 pin 到第 $i\bmod m$ 个核，用于观察是否被重新收敛。
\end{itemize}

配套的 \texttt{yat\_affinity\_test.c} 在 guest 内直接验证"返回的核 $\subseteq$ 允许掩码"，输出 \texttt{YAT-AFFINITY PASS/FAIL} 与三个阶段的实际核编号，作为亲和性修复的回归测试。
